% !PRINTED 126
\weberpage{126}

%<<CONT
sein. Denn setzen wir $t = m s + s_1$, worin $s_1 < s$ ist, so ist $c^{s_1}$
in $Q$ enthalten, und $s_1$ muss also $= 0$ sein. Demnach ist $s$ ein
Theiler von $\gamma$. Nun bilden aber die Elemente

\[
Q + c\,Q + c^2\,Q + \cdots c^{s-1}\,Q
\]

wegen (1) eine in $P$ enthaltene Gruppe, und zwar vom Grade $p^{\varkappa} s$,
und da also $p^{\varkappa} s$ ein Theiler von $n$ sein muss, so folgt, dass $s$
nicht durch $p$ theilbar sein kann. Wäre nun der Grad $\gamma$ von $c$
eine Potenz von $p$, so müsste auch $s$ eine Potenz von $p$ sein, was
hiernach nicht möglich ist.

Es kann also der Grad von $c$ gewiss nicht eine Potenz von
$p$ sein; d. h. \gesperrt{ausser den Elementen $Q$ giebt es in $R$ keine
Elemente, deren Grad eine Potenz von $p$ ist.}

Bedeutet nun $b$ ein Element von $P$, das nicht zu $R$ gehört,
so giebt es nach der Definition von $R$ in der Gruppe

\[
Q' = b^{-1}\,Q\,b
\]

mindestens ein Element, das nicht in $Q$ vorkommt, und der
Grad dieses Elementes muss eine Potenz von $p$ sein, weil der
Grad von $Q'$ gleich $p^{\varkappa}$, also eine Potenz von $p$ ist. Dies Element
kann also nicht in $R$ vorkommen, \gesperrt{und $Q'$ ist daher kein
Theiler von $R$.}

Wählen wir die Elemente $b_1, b_2, \ldots, b_{j-1}$ aus $P$ so aus,
dass wir

\weq{10}{P = R + R\,b_1 + R\,b_2 + \cdots + R\,b_{j-1}} % [sic: Weber jumps from equation (2) to (10) in this section and subsequently cites (10).]

setzen können, so sind also die $j - 1$ Gruppen:

\[
b_1^{-1}\,Q\,b_1,\ \ b_2^{-1}\,Q\,b_2,\ \ \ldots,\ \ b_{j-1}^{-1}\,Q\,b_{j-1},
\]

die alle vom Grade $p^{\varkappa}$ sind, von $Q$ verschieden. Sie sind aber
auch von einander verschieden; denn wäre z. B.

\[
b_1^{-1}\,Q\,b_1 = b_2^{-1}\,Q\,b_2,
\]

so würde folgen:

\[
b_2\,b_1^{-1}\,Q\,b_1\,b_2^{-1} = Q,
\]

d. h. $b_2\,b_1^{-1} = c$ wäre in $R$ enthalten, also $b_2 = c\,b_1$ in $R\,b_1$, was
wegen (10) nicht möglich ist. Wir haben also gewiss $j$ verschiedene
conjugirte Gruppen $p^{\varkappa\,\mathrm{ten}}$ Grades:

\weq{11}{Q,\ \ b_1^{-1}\,Q\,b_1,\ \ b_2^{-1}\,Q\,b_2,\ \ \ldots,\ \ b_{j-1}^{-1}\,Q\,b_{j-1}.}

Ist $Q' = b^{-1}\,Q\,b$ eine von ihnen, so enthält $R' = b^{-1}\,R\,b$
den Theiler $Q'$, aber ausser diesem kein Element, dessen Grad
%CONT>>
